
All levels share the objective function and component models set out below. They differ
only in two places (the topology of the energy balance and the treatment of thermal
losses), and it is precisely because these two enter as separable terms that the
decomposition of Section~\ref{sec:results} is exact rather than fitted.

\paragraph{Objective function.}
The optimiser minimises total annual operating cost: electricity procurement net of
feed-in revenue, fuel, heat-curtailment penalties, carbon charges, a demand charge on
peak grid import, and a thermal-storage cycling penalty,
\begin{equation}
  \min \; Z = C^{\mathrm{en}} + C^{\mathrm{fuel}}
    + C^{\mathrm{dump}} + C^{\mathrm{CO_2}} + C^{\mathrm{dem}}
    + C^{\mathrm{cyc}} .
  \label{eq:objective}
\end{equation}
Network pumping does not appear as a separate cost term. The pump electrical power
$P_t^{\mathrm{pump}}$ enters the electricity balance, Eq.~\eqref{eq:balance_el}, and is
therefore paid for through grid procurement inside $C^{\mathrm{en}}$; it is identically
zero at every level below $\Lthree$ and is computed from Eq.~\eqref{eq:pump_power} from
$\Lthree$ upward (Section~\ref{sec:extended}). With $\Delta t = 1$\,h and $t$ indexing all
8\,760 hours of the year,
\begin{align*}
C^{\mathrm{en}}   &= \Delta t \sum_t
  \bigl(P_t^{\mathrm{buy}} \lambda_t^{\mathrm{buy}}
        - P_t^{\mathrm{sell}} \lambda_t^{\mathrm{sell}}\bigr), &
C^{\mathrm{fuel}} &= \Delta t \sum_{g,t}
  F_{g,t}\, \lambda_g^{\mathrm{fuel}}, \\
C^{\mathrm{dump}} &= \Delta t \sum_{n,t}
  Q_{n,t}^{\mathrm{dump}}\, \lambda^{\mathrm{dump}}, &
C^{\mathrm{dem}}  &= \lambda^{\mathrm{dem}} f^{\mathrm{yr}}
  P^{\mathrm{peak}}, \\
C^{\mathrm{CO_2}} &= \frac{\lambda^{\mathrm{CO_2}}}{1000} \sum_t
  \Bigl(E_t^{\mathrm{CO_2,grid}} +
        \textstyle\sum_g E_{g,t}^{\mathrm{CO_2,fuel}}\Bigr), &
C^{\mathrm{cyc}} &= \lambda^{\mathrm{cyc}}\,\Delta t \sum_{s,t}
  \bigl(Q_{s,t}^{\mathrm{ch}} + Q_{s,t}^{\mathrm{dis}}\bigr).
\end{align*}
Here $P_t^{\mathrm{buy}}$ and $P_t^{\mathrm{sell}}$ [MW] are grid import and export at
time $t$ and $\lambda_t^{\mathrm{buy}}$, $\lambda_t^{\mathrm{sell}}$ [EUR/MWh] the
corresponding prices; $F_{g,t}$ [MW] is the fuel input to generator $g$ on a
heating-value basis at specific cost $\lambda_g^{\mathrm{fuel}}$;
$Q_{n,t}^{\mathrm{dump}}$ [MW] is heat curtailed at node $n$, penalised at
$\lambda^{\mathrm{dump}}$; $\lambda^{\mathrm{CO_2}}$ [EUR/t] is the carbon price;
$f^{\mathrm{yr}}$ an annualisation factor; $P^{\mathrm{peak}}$ [MW] the annual peak
grid import entering the demand charge $\lambda^{\mathrm{dem}}$; and
$\lambda^{\mathrm{cyc}}$ [EUR/MWh] the storage-cycling cost, applied to the combined
charge-and-discharge throughput $Q_{s,t}^{\mathrm{ch}}+Q_{s,t}^{\mathrm{dis}}$ of each
storage $s$.

Two elements of $Z$ depart from operator cash flow: the carbon charge $C^{\mathrm{CO_2}}$
enters on a gross basis, and the cycling penalty $C^{\mathrm{cyc}}$ is a wear-and-arbitrage
deterrent rather than a cash outlay. The reported \emph{economic} cost applies the standard
cogeneration self-use carbon credit and excludes $C^{\mathrm{cyc}}$, so it differs
systematically from $Z$; the economic cost is used for every reported comparison. The
distinction and its magnitude are set out in Section~\ref{subsec:costacct}.

\paragraph{Energy balance, and how the controls are constructed.}
The copperplate levels $\CP$ and $\CPL$ enforce a single bus-level balance over all
generators, heat pumps, electrode boilers and storage, with no spatial routing:
\begin{equation}
  \sum_{g} Q_{g,t}^{\mathrm{th}}
    + \sum_{h} Q_{h,t} + \sum_{r} Q_{r,t}
    + \sum_{s} Q_{s,t}^{\mathrm{dis}}
  = Q_t^{\mathrm{dem}} + Q_t^{\mathrm{dump}}
    + \sum_{s} Q_{s,t}^{\mathrm{ch}}
  + \underbrace{L_t}_{\CPL\ \mathrm{only}}
  \;\; \forall\, t,
  \label{eq:balance_cp}
\end{equation}
where $Q_{g,t}^{\mathrm{th}}$ [MW] is the thermal output of generator $g$,
$Q_{h,t}$ and $Q_{r,t}$ [MW] the heat-pump and electrode-boiler thermal outputs,
$Q_{s,t}^{\mathrm{dis}}$ and $Q_{s,t}^{\mathrm{ch}}$ [MW] the discharge and charge of
storage $s$, and $Q_t^{\mathrm{dem}}$ [MW] the total network demand. The term $L_t$ is
the exogenous aggregate loss supplied to the control condition $\CPL$ and is zero for
$\CP$; its calibration is given in Section~\ref{subsec:lumped}.

All node-resolved levels ($\ZN$, $\NDz$ and $\Lone$ through $\Lsix$) instead enforce
a nodal balance at each junction $n$, accounting for thermal losses along incoming pipes
and for the spatial distribution of assets and demands:
\begin{align}
  &\sum_{p \in \mathcal{P}_n^{\mathrm{in}}}
    \bigl(Q_{p,t} - \kappa\, Q_{p,t}^{\mathrm{loss,pipe}}\bigr)
  + \sum_{g \in \mathcal{G}_n} Q_{g,t}^{\mathrm{th}}
  + \sum_{h \in \mathcal{H}_n} Q_{h,t}
  + \sum_{r \in \mathcal{R}_n} Q_{r,t}
  + \sum_{s \in \mathcal{S}_n} Q_{s,t}^{\mathrm{dis}} \notag \\
  &= Q_{n,t}^{\mathrm{dem}} + Q_{n,t}^{\mathrm{dump}}
  + \sum_{s \in \mathcal{S}_n} Q_{s,t}^{\mathrm{ch}}
  + \sum_{p \in \mathcal{P}_n^{\mathrm{out}}} Q_{p,t}
  \;\; \forall\, n,\, t,
  \label{eq:balance_nodal}
\end{align}
with $\mathcal{P}_n^{\mathrm{in}}$, $\mathcal{P}_n^{\mathrm{out}}$ the pipes entering and
leaving node $n$, $Q_{p,t}$ [MW] the heat flow in pipe $p$, and
$Q_{p,t}^{\mathrm{loss,pipe}}$ [MW] its thermal loss.

The switch $\kappa \in \{0,1\}$ is the point of the design. Equations
\eqref{eq:balance_cp} and \eqref{eq:balance_nodal} differ in two respects at once
(spatial resolution and loss visibility), and prior fidelity comparisons have moved both
together, which is why their cost differences cannot be attributed to either. Setting the
two independently generates the four corners of a $2\times2$ design: $\CP$ is
\eqref{eq:balance_cp} with $L_t = 0$ (neither), $\CPL$ is \eqref{eq:balance_cp} with
$L_t$ supplied (loss, no topology), $\NDz$ is \eqref{eq:balance_nodal} with $\kappa = 0$
(topology, no loss), and $\Lone$ is \eqref{eq:balance_nodal} with $\kappa = 1$ (both).
The cost gap $\CP \to \Lone$ therefore decomposes additively into a loss main effect, a
topology main effect and their interaction, with no residual by construction. This is the
control condition the confound requires, and it is the reason the decomposition in
Section~\ref{subsec:decomp} closes to machine precision rather than to a fitted
tolerance.

The global electricity balance and grid mutual-exclusivity constraints close the system.
All electrical generators, power-to-heat units $r$, heat pumps and pumps balance against
grid import, export and in-network consumption:
\begin{align}
  P_t^{\mathrm{buy}}
    &+ \sum_{g \in \mathcal{G}^{\mathrm{CHP}}} P_{g,t}^{\mathrm{el}}
  = \sum_{h \in \mathcal{H}} P_{h,t}^{\mathrm{el}}
    + \sum_{r \in \mathcal{R}} P_{r,t}^{\mathrm{el}} \notag \\
    &+ P_t^{\mathrm{pump}} + P_t^{\mathrm{sell}}
  \quad \forall\, t,
  \label{eq:balance_el} \\[4pt]
  P_t^{\mathrm{buy}} &\leq M^{\mathrm{grid}}\, z_t,
  \;\;
  P_t^{\mathrm{sell}} \leq M^{\mathrm{grid}}\, (1 - z_t)
  \;\; \forall\, t,
  \label{eq:grid_bigm}
\end{align}
where $P_{g,t}^{\mathrm{el}}$ [MW] is CHP electrical output, $P_{h,t}^{\mathrm{el}}$ the
heat-pump electrical input, $P_{r,t}^{\mathrm{el}}$ the electrode-boiler input and
$P_t^{\mathrm{pump}}$ the network pumping power. The binary $z_t \in \{0,1\}$ enforces
that the grid connection either imports or exports in any hour, with $M^{\mathrm{grid}}$
a Big-M set to the grid capacity.

\paragraph{Thermal loss model.}
For each pipe $p$, steady-state conductive losses are computed separately for the supply
and return circuit. Because the supply temperature follows the heating curve rather than
being optimised, the losses are fixed parameters, the simplification that keeps the
formulation linear:
\begin{align}
  Q_{p,t}^{\mathrm{loss,sup}} &= \frac{U^{s}_p\, L_p\,
    \bigl(T_{\mathrm{sup}}^{\mathrm{nom}}(t)
    - T_{\mathrm{ground}}(t)\bigr)}{10^6},
  \label{eq:loss_supply} \\[4pt]
  Q_{p,t}^{\mathrm{loss,ret}} &= \frac{U^{r}_p\, L_p\,
    \bigl(T_{\mathrm{ret}}^{\mathrm{nom}}(t)
    - T_{\mathrm{ground}}(t)\bigr)}{10^6},
  \label{eq:loss_return}
\end{align}
with $U^{s}_p$, $U^{r}_p$ [W/(m\,K)] the linear heat-transfer coefficients of the supply
and return pipe, set by diameter and insulation class, $L_p$ [m] the length,
$T_{\mathrm{sup}}^{\mathrm{nom}}$, $T_{\mathrm{ret}}^{\mathrm{nom}}$ [°C] the nominal
temperatures prescribed by the heating curve and $T_{\mathrm{ground}}$ [°C] the ground
temperature; $10^6$ converts W to MW. Losses are identically zero for $\CP$ and $\NDz$,
enter exogenously as $L_t$ for $\CPL$, and are computed by \eqref{eq:loss_supply} and
\eqref{eq:loss_return} for $\ZN$ and $\Lone$ upward.

Two points about the calibration matter for what follows. First, the $U$ values are
constrained to a defensible range for the pipe class rather than fitted freely: forcing the
trunk-only loss to match the total inferred from the delivered-energy record by inflating
terminal-pipe coefficients cannot be attributed to insulation ageing. Second, what such a multiplier would absorb is the
\emph{last mile}, the service laterals and substations a trunk-only model does not represent,
and it is therefore placed where it physically belongs, entering only at the station-resolved
level $\Lfour$. The node-resolved levels consequently \emph{undercount} the network
loss by construction, and the modelled total matches the measured \emph{delivered energy} to
about one percent only once the laterals are added (Section~\ref{subsec:validation}); the
network loss this implies is model-derived, not independently metered. This is a result rather
than a defect: a quantitative statement of what a trunk-resolved model cannot see.

\paragraph{Heat pump.}
Each heat pump $h$ either recovers industrial waste heat at a time-varying source
temperature or falls back on a low-temperature source at constant performance, the
thermal output being the sum of the two contributions:
\begin{align}
  Q_{h,t} &= Q_{h,t}^{\mathrm{WHR}} + Q_{h,t}^{\mathrm{def}},
  \label{eq:hp_split} \\[4pt]
  P_{h,t}^{\mathrm{el}} &= \frac{Q_{h,t}^{\mathrm{WHR}}}{
    \mathrm{COP}_{h,t}^{\mathrm{WHR}}}
    + \frac{Q_{h,t}^{\mathrm{def}}}{\mathrm{COP}_h^{\mathrm{def}}}.
  \label{eq:hp_pel}
\end{align}
The waste-heat coefficient of performance $\mathrm{COP}_{h,t}^{\mathrm{WHR}}$ is
pre-computed from an analytical Lorenz-cycle model
(Appendix~\ref{app:cop}), which avoids bilinear terms in the optimisation. That
pre-computation is a deliberate scope restriction rather than an approximation of
convenience, and it is revisited in the limitations: a model with endogenous COP would
give the heat pump more room to respond to the network state than it has here.

Operation is bounded by a minimum load fraction $\phi_h^{\mathrm{min}}$ and the installed
capacity $\bar{Q}_h$, switched by the on/off binary $u_{h,t} \in \{0,1\}$:
\begin{equation}
  \phi_h^{\mathrm{min}}\, \bar{Q}_h\, u_{h,t}
  \leq Q_{h,t} \leq \bar{Q}_h\, u_{h,t}
  \quad \forall\, h,\, t,
  \label{eq:hp_cap}
\end{equation}
so that the unit is either off or above its minimum load, preventing uneconomic part-load
cycling.

\paragraph{Thermal storage.}
The state of charge $E_{s,t}$ [MWh] of storage unit $s$ follows a discrete-time balance
with self-discharge and separate charging and discharging efficiencies:
\begin{equation}
  E_{s,t} = (1 - \alpha_s)^{\Delta t/\mathrm{1\,h}}\, E_{s,t-1}
    + \eta_s^{\mathrm{ch}}\, Q_{s,t}^{\mathrm{ch}}\, \Delta t
    - \frac{Q_{s,t}^{\mathrm{dis}}\, \Delta t}{\eta_s^{\mathrm{dis}}}.
  \label{eq:storage_soc}
\end{equation}
The coefficient $\alpha_s$, the fractional standing loss per hour, represents losses of
the vessel, and the exponential form makes the discretisation independent of the time step.
A cyclic condition $E_{s,|\mathcal{T}|} = E_{s,0}$ enforces energy-neutral operation over
the year, preventing artificial depletion or accumulation.

\paragraph{Remaining assets and emissions.}
Gas boiler, biomass boiler, electrode boiler and CHP follow standard linear formulations,
capacity-bounded with constant efficiency, detailed in
Appendix~\ref{app:component_models}. Grid and fuel emissions are the product of energy
consumed and the corresponding emission intensity,
\begin{equation}
  E_t^{\mathrm{CO_2,grid}} =
    P_t^{\mathrm{buy}}\, e_t^{\mathrm{grid}}\, \Delta t,
  \qquad
  E_{g,t}^{\mathrm{CO_2,fuel}} =
    F_{g,t}\, e_g^{\mathrm{fuel}}\, \Delta t,
  \label{eq:emissions}
\end{equation}
with $e_t^{\mathrm{grid}}$ [kg/MWh] the hourly marginal grid emission factor and
$e_g^{\mathrm{fuel}}$ [kg/MWh] the fuel-specific intensity. Both are fixed inputs: the
optimiser is carbon-aware but does not optimise the emission factor itself.
